\section{Related Work}

\subsection{Data Collection Paradigms for Embodied Intelligence}

The pursuit of generalizable embodied intelligence hinges on the availability of large-scale, diverse, and high-quality demonstration data~\cite{lindata,gen0}. Current data collection paradigms can be broadly categorized into four approaches. \textbf{Robot teleoperation}~\cite{chen2025towards} yields high-fidelity demonstrations but is constrained by expensive hardware and lab environments, with efforts like DROID~\cite{droid2024} pooling 76K trajectories across 350 hours. \textbf{Universal Manipulation Interfaces (UMIs)}~\cite{chi2024umi,xu2025dexumi} reduce costs via portable grippers but remain ``active'' devices requiring deliberate manipulation, limiting natural interaction coverage. \textbf{Wearable AR/VR systems}~\cite{wang2024dexcap,zhong2025humanoidexo,fang2024airexo} capture complex hand poses with high precision, yet their significant bulk ($>500$g) and power dependencies render them too intrusive for continuous daily use. \textbf{Passive egocentric video} from existing datasets~\cite{grauman2022ego4d,damen2018epic} offers massive scale but typically lacks fine-grained manipulation annotations, imposing prohibitive curation costs for robot training.

Our AoE system~\cite{aoe2026} introduced a paradigm shift by leveraging ubiquitous smartphones for passive, always-on egocentric data collection at near-zero marginal cost ($<\$20$ per user), combining the scalability of passive video capture with the annotation richness required for policy learning.

\subsection{Egocentric Datasets for Manipulation Learning}

The landscape of egocentric manipulation datasets has expanded rapidly. \textbf{Generalist datasets} such as Ego4D~\cite{grauman2022ego4d} (3,670 hours across 74 locations) and EPIC-KITCHENS~\cite{damen2018epic} (100 hours) provide broad coverage but lack manipulation-specific geometric annotations. \textbf{Manipulation-centric datasets} have emerged to fill this gap: EgoDex~\cite{egodex2025} collected 829 hours via Apple Vision Pro with native 3D hand tracking across 194 tasks; EgoLive~\cite{egolive2026} offers 1,680 hours of stereo video at 60 FPS from real-world service scenarios; and OpenEgo~\cite{openegocentric2025} unifies six public datasets (1,107 hours) with standardized 21-joint hand poses. At the \textbf{deployment scale}, NVIDIA's EgoScale~\cite{egoscale2026} demonstrates log-linear scaling behavior with over 20,000 hours.

A parallel thread focuses on learning from egocentric data for robot control. Being-H0~\cite{luo2025being} pioneers VLA pretraining via MANO-based motion tokenization. EgoVLA~\cite{egovla2025} and VITRA~\cite{vitra2025} construct robot-aligned training data from egocentric videos. Despite these advances, most existing datasets share two limitations: reliance on specialized, expensive capture devices, and absence of toolchains bridging raw data to model training. Open-AoE addresses both.

\subsection{Human-to-Robot Transfer and Cross-Embodiment Learning}

Bridging the embodiment gap between human demonstrations and robot execution spans kinematic retargeting~\cite{gmr2026,dexretarget}, physics-based simulation~\cite{spider2025,maniptrans2025}, and end-to-end policy transfer~\cite{egomimic2025,twist2025}. For mainstream VLA model training, the community has converged on several de facto standards: $\pi_{0.5}$~\cite{pi05} and GR00T~\cite{groot2025} for foundation model architectures, LeRobot~\cite{lerobot2024} for data format standardization, and frameworks like FLARE~\cite{flare2025} for cross-embodiment fine-tuning. Open-AoE's toolchain directly targets compatibility with these ecosystems.
